\begin{table}[t]
\centering
\caption{0.5\,ha 컷오프에서 사전결정 공변량의 연속성 검정 (CCT 2014 \texttt{rdrobust}, 삼각 커널, 클러스터 = \texttt{hh\_id}). 2018년 기준연도 단면.}
\label{tab:cct_covariate_continuity_ko}
\begin{tabular}{lrrr}
\toprule
Bandwidth & 계수 (robust) & SE & $N$ \\
\midrule
\multicolumn{4}{l}{\textit{세대주 연령 등급}} \\
T1 ($h=500$) & 0.299 & (0.379) & 142 \\
T2 ($h=1,000$) & 0.266 & (0.315) & 293 \\
T3 ($h=1,627$) & 0.145 & (0.210) & 480 \\
\multicolumn{4}{l}{\textit{영농 유형}} \\
T1 ($h=500$) & 0.946 & (1.706) & 142 \\
T2 ($h=1,000$) & -0.681 & (1.228) & 293 \\
T3 ($h=1,523$) & -1.479* & (0.822) & 449 \\
\multicolumn{4}{l}{\textit{전\textperiodcentered{}겸업 구분}} \\
T1 ($h=500$) & 0.057 & (0.486) & 142 \\
T2 ($h=1,000$) & -0.269 & (0.328) & 293 \\
T3 ($h=1,554$) & -0.352* & (0.211) & 456 \\
\multicolumn{4}{l}{\textit{기준연도 자작 비율 ($s_{0,i}$)}} \\
T1 ($h=500$) & 0.060 & (0.166) & 142 \\
T2 ($h=1,000$) & 0.130 & (0.135) & 293 \\
T3 ($h=1,613$) & -0.011 & (0.092) & 476 \\
\multicolumn{4}{l}{\textit{기준연도 자작 면적}} \\
T1 ($h=500$) & 1,142.522 & (1,584.988) & 142 \\
T2 ($h=1,000$) & 1,248.293 & (1,101.498) & 293 \\
T3 ($h=1,522$) & -441.111 & (739.728) & 449 \\
\multicolumn{4}{l}{\textit{기준연도 농외소득}} \\
T1 ($h=500$) & -1,941,823.253 & (6,372,928.771) & 142 \\
T2 ($h=1,000$) & -5,290,487.756 & (4,371,873.648) & 293 \\
T3 ($h=1,737$) & -4,160,397.784 & (2,566,226.301) & 518 \\
\bottomrule
\end{tabular}
\\
\footnotesize\textit{주: 각 사전결정 공변량을 $rv_{2018,i} = 0$에서 평가한 robust 편향보정 계수 (\citealp{CalonicoCattaneoTitiunik2014_rdrobust}). $N$은 선택된 bandwidth 내 유효 표본 (컷오프 좌우 합). T3 = 결과변수별 MSE 최적. 클러스터-강건 SE는 \texttt{hh\_id} 기준. * p$<$0.10, ** p$<$0.05, *** p$<$0.01.}
\end{table}
